\documentclass[11pt,a4paper]{article}
\usepackage[T1]{fontenc}
\usepackage{lmodern,amsmath,amssymb,amsthm,mathtools,microtype}
\usepackage[margin=25mm]{geometry}
\usepackage{xcolor}
\usepackage[colorlinks=true,linkcolor=blue!45!black,citecolor=blue!45!black,urlcolor=blue!45!black]{hyperref}
\newtheorem{theorem}{Theorem}[section]
\newtheorem{proposition}[theorem]{Proposition}
\newtheorem{lemma}[theorem]{Lemma}
\newtheorem{corollary}[theorem]{Corollary}
\theoremstyle{definition}\newtheorem{definition}[theorem]{Definition}
\newcommand{\Q}{\mathbb Q}\newcommand{\Z}{\mathbb Z}
\setlength{\emergencystretch}{2em}
\hypersetup{pdftitle={Uniform rational approximations to an elliptic logarithm}}
\title{Uniform Rational Approximations\\to an Elliptic Logarithm}
\author{}\date{8 November 2026}
\begin{document}\maketitle
\begin{abstract}
For the elliptic curve $y^2=x^3-Dx$, we give an explicit rational
polynomial approximation to the formal logarithm, with a local error
bound uniform in the truncation length. The estimate applies at every
odd unramified prime and does not require the prime to exceed the
truncation order. The same polynomials are positive on the real
interval $[0,1]$, which proves nonvanishing of the associated
logarithms at every nonzero double on the specified curve over
$\mathbb Q(\sqrt2)$. Integer coefficient formulas and norm bounds
provide exact finite specifications for arbitrarily long truncations.
All statements in this paper are unconditional.
\end{abstract}
\section{Uniform logarithm truncations of arbitrary order}
\label{sec:uniform-logarithm}

Let $E:y^2=x^3-Dx$ over a field of characteristic zero, with $D\ne0$.
Write $f=1/x$ and $z=-x/y$ near the identity. The equation gives
\[
                         f=z^2(1-Df^2).
\]
Define rational numbers
\[
 c_0=1,\qquad c_1=-\frac25,\qquad
 c_j=c_{j-1}\frac{4j-5}{4j+1}\quad(j\ge2),
\]
and put
\[
 Q_J(T)=\sum_{j=0}^{J-1}c_jT^j,
 \qquad L_J(Q)=z(Q)Q_J(Df(Q)^2)\quad(J\ge1).
\]

\begin{lemma}\label{lem:logarithm-series}
The formal series
\[
                  L(Q)=z(Q)\sum_{j\ge0}c_j(Df(Q)^2)^j
\]
is the normalized formal logarithm of the displayed equation.
For every odd prime $p$ and $j\ge1$,
\[
                  v_p(c_j)\ge-\lfloor\log_p(4j+1)\rfloor.
\]
\end{lemma}
\begin{proof}
The invariant differential is
$\omega=dx/(2y)=z\,df/(2f)$. Put $T=Df^2$ and $S(T)=\sum c_jT^j$.
From $z^2=f/(1-T)$, differentiation shows that $d(zS(T))=\omega$
is equivalent to
\[
                 4T(1-T)S'(T)+(1+T)S(T)=1-T.
\]
The constant coefficient is $c_0=1$, the linear coefficient is
$5c_1=-2$, and every subsequent coefficient is precisely the stated
recurrence. The resulting series is $z+O(z^5)$ and hence is the
normalized formal logarithm.

For $j\ge1$ its coefficient is exactly
\[
                 c_j=2\prod_{i=1}^j\frac{4i-5}{4i+1}.
\]
For an odd prime power $p^r$, each of the two arithmetic progressions
in this quotient has length $j$ and invertible step four. The counts
of terms divisible by $p^r$ therefore differ by at most one.
No term has absolute value exceeding $4j+1$. Summing the possible
loss of one valuation over the prime powers at most $4j+1$ proves
the claimed lower bound.
\end{proof}

\begin{theorem}\label{thm:uniform-logarithm-tail}
Let $F/\Q_p$ be unramified, where $p$ is odd, normalize $v(p)=1$,
and suppose $D$ is integral. For any nonzero point $Q$ in the formal
group, so that $v(z(Q))\ge1$, and every integer $J\ge1$,
\[
 \begin{aligned}
 v(L_J(Q))&=v(z(Q)),\\
 v(L(Q)-L_J(Q))&\ge v(z(Q))+3J.
 \end{aligned}
\]
In particular, the conclusion is uniform in $J$ and does not require
$p>J$ or $p>4J$.
\end{theorem}
\begin{proof}
The relation between $f$ and $z$ gives $v(f(Q))=2v(z(Q))$.
For every $j\ge1$ and odd $p$ one has
\[
                 \lfloor\log_p(4j+1)\rfloor\le j.
\]
Indeed $3^{j+1}>4j+1$ for $j\ge1$, by induction. Consequently
Lemma~\ref{lem:logarithm-series} gives
\[
 \begin{split}
 v\bigl(z(Q)c_j(Df(Q)^2)^j\bigr)
 &\ge v(z(Q))+4jv(z(Q))-\lfloor\log_p(4j+1)\rfloor\\
 &\ge v(z(Q))+3j.
 \end{split}
\]
All nonconstant terms are strictly smaller in absolute value than
the leading term $z(Q)$, proving the first equality. The same estimates
show that the infinite series converges, and that its tail beginning
at $j=J$ has the stated valuation.
\end{proof}

\begin{corollary}\label{cor:uniform-logarithm-ratio}
In the setting of Theorem~\ref{thm:uniform-logarithm-tail}, suppose
$Q=[n]P$ and $R=[m]P$ are nonzero formal points, where $m,n$ are
nonzero integers. For truncation lengths $J_1,J_2\ge1$, put
$k=3\min(J_1,J_2)$. Then
\[
 v\left(\frac{L_{J_1}(Q)}{L_{J_2}(R)}-\frac nm\right)
                  \ge v(n/m)+k.
\]
If $n/m$ is integral, the right side is at least $k$.
\end{corollary}
\begin{proof}
Formal additivity applied to $[m]Q=[n]R$ gives $L(Q)/L(R)=n/m$.
The leading-term estimate in the theorem makes both logarithms
nonzero. Write $L_{J_1}(Q)=L(Q)(1+\varepsilon_Q)$ and
$L_{J_2}(R)=L(R)(1+\varepsilon_R)$; their two errors have valuations
at least $k$. Subtracting the exact ratio proves the assertion.
\end{proof}

For the field $\Q(\sqrt2)$, every odd completion is unramified and
$D=11-6\sqrt2$ is integral. Thus these statements apply at every
contact prime $p>76176$, independently of the chosen truncation
length. If $M$ is a power of two with $M\ge2$, the choice $J=M/2$
has relative precision at least $3M/2$, which is stronger than
precision $M$. Arbitrarily long truncations therefore leave the fixed
contact-height constant and its exceptional-prime set unchanged.

\subsection{Nonvanishing at the real embeddings}

\begin{lemma}\label{lem:positive-logarithm-polynomial}
For every integer $J\ge1$, one has $Q_J(T)>0$ for $0\le T\le1$.
\end{lemma}
\begin{proof}
All $c_j$ for $j\ge1$ are negative. For $r\ge1$, the recurrence
telescopes to
\[
                    \sum_{j=0}^{r}c_j=-\frac{4r-1}{2}c_r>0.
\]
The identity holds at $r=1$ and its successive differences are $c_r$.
Because $0\le T^j\le1$ and $c_j<0$, it follows that
$Q_J(T)\ge Q_J(1)>0$; $J=1$ is immediate.
\end{proof}

\begin{corollary}\label{cor:nonzero-logarithm-denominator}
For $K=\Q(\sqrt2)$ and $D=(3-\sqrt2)^2$, let $Q\in2E(K)$
be nonzero. Then $L_J(Q)\ne0$ for every $J\ge1$.
\end{corollary}
\begin{proof}
The duplication identity
\[
                x(2R)=\left(\frac{x(R)^2+D}{2y(R)}\right)^2
\]
shows that the abscissa of a nonzero double is positive at both real
embeddings. Its ordinate cannot be zero: the two nonzero
ordinate-zero abscissas are $\pm(3-\sqrt2)$, neither a square in $K$
since their norm is seven, while abscissa zero is excluded by the
duplication formula. At either embedding the curve equation now
implies $x(Q)>\sqrt D$. Thus $0<Df(Q)^2<1$ and $z(Q)\ne0$.
Lemma~\ref{lem:positive-logarithm-polynomial} proves the assertion.
\end{proof}

Consequently any rational system whose logarithm denominators are
evaluated at nonzero doubles may use these truncations at arbitrary
finite length. Increasing the length introduces no new
nonvanishing condition at those points.

\subsection{A finite integer coefficient specification}

For a polynomial $G$, write $\|G\|_1$ for the sum of the absolute
values of its coefficients. In a finite free algebra over $\mathbb Q$
or an integer polynomial ring, choose a basis and let $M_A$ be the
matrix of multiplication by $A$. Use the maximum column sum of the
coefficient $\ell^1$ norms as its matrix norm. This norm is
submultiplicative, has $\|I\|=1$, and satisfies $\|cI\|=|c|$ for
rational scalars.

Compatible weights mean nonnegative weights on base variables and
positive weights on the algebra generators for which multiplication
and reduction by the monic defining relations do not increase the
sum of the input weights. The weighted degree of the zero element is
$-\infty$; all other weighted degrees are nonnegative.

Define
\[
 D_J=\prod_{i=1}^{J-1}(4i+1),\qquad
 G_J(T)=D_J Q_J(T)\in\Z[T].
\]
Empty products are one. Its coefficients are explicitly
\[
 [T^0]G_J=D_J,\qquad
 [T^j]G_J=
 2\left(\prod_{i=1}^{j}(4i-5)\right)
  \left(\prod_{i=j+1}^{J-1}(4i+1)\right)
 \quad(1\le j<J).
\]
This finite product-and-sum definition specifies the ordinary integer
polynomial without requiring a list of all its coefficients.

\begin{proposition}\label{prop:indexed-logarithm-coefficients}
For $J\ge2$, the degree of $G_J$ is $J-1$ and
\[
 D_J\le(4J)^{J-1},\qquad
              \|G_J\|_1\le2J(4J)^{J-1}.
\]
If $J=2^{b-1}$, with $b\ge2$, these imply
\[
                    D_J<2^{(b+1)J},\qquad
                    \|G_J\|_1<2^{(b+1)J}.
\]
If $A$ is an element of a finite free algebra with
$\|M_A\|\le2^B$, $B\ge0$, then
\[
                  \|M_{G_J(A)}\|
                  \le\|G_J\|_1\,2^{(J-1)B}.
\]
If the algebra has the multiplicative weighted filtration associated
with compatible positive generator weights, then
$\operatorname{wdeg}(G_J(A))\le(J-1)\max(0,\operatorname{wdeg}(A))$.
\end{proposition}
\begin{proof}
Every nonconstant coefficient is a product of $J-1$ integers of
absolute value at most $4J$, times two. The same bound without two
holds for the constant coefficient. Summing the coefficient bounds
gives the displayed estimate. Its binary exponent when
$J=2^{b-1}$ is
\[
       1+(b-1)+(J-1)(b+1)=(b+1)J-1.
\]
The operator estimate follows from the triangle inequality and
submultiplicativity, and the weight estimate follows term by term.
\end{proof}

For a rational argument $A/B$, its homogeneous numerator is the
explicit polynomial
\[
 G_J^{\rm hom}(A,B)=\sum_{j=0}^{J-1}[T^j]G_J\,A^jB^{J-1-j}.
\]
Thus $Q_J(A/B)=G_J^{\rm hom}(A,B)/(D_JB^{J-1})$ on $B\ne0$.
This form is convenient for a finite algebra over an integer polynomial
ring. Its degree and operator-norm estimates follow from the same
proof with the larger of the two argument bounds.

\end{document}
